\documentclass[11pt]{article}

\usepackage[margin=0.85in]{geometry}
\usepackage{booktabs}
\usepackage{array}
\usepackage{graphicx}
\usepackage{hyperref}
\usepackage{xcolor}

\newcommand{\barrow}[4]{#1 & #2 & \textcolor{#4}{\rule{#3cm}{7pt}}\\}

\title{SPY TP2/RR2 Relative-Value Pipeline Summary}
\author{FinalQuant}
\date{May 2026}

\begin{document}
\maketitle

\section{Scope}

This report summarizes the full end-of-day SPY option pipeline that was implemented and run on the local archive:

\begin{center}
\texttt{archive/spy\_options\_data\_19.json} through \texttt{archive/spy\_options\_data\_24.json}
\end{center}

The universe is SPY ETF options, end-of-day frequency, years 2019--2024, DTE 1--60, excluding 0DTE. Detection uses the two-strike TP2/RR2 approximation. This is SPY ETF option relative-value research, not pure SPX arbitrage replication. SPY options are American-style and affected by ETF, dividend, spread, and close-time quote effects.

\section{What Was Built}

\begin{itemize}
  \item Ingested and flattened yearly SPY option JSON archives.
  \item Normalized quote fields, Greeks, volume, open interest, bid/ask sizes, mid, spread, spread percent, and DTE.
  \item Applied base quote filters: valid bid/ask, positive mid, DTE 1--60, non-extreme spreads, and reasonable absolute delta.
  \item Attached official raw SPY close when available, with put-call parity proxy fallback for missing dates.
  \item Built daily call/put chains grouped by expiration and aligned strike pairs.
  \item Detected TP2 call violations and RR2 put violations using mid prices.
  \item Generated violation features, leg features, liquidity aggregates, IV/Greek slopes, and persistence features.
  \item Expanded each violation into five strategy rows: \texttt{full4}, \texttt{T1}, \texttt{T2}, \texttt{K1}, and \texttt{K2}.
  \item Computed 5-trading-day after-cost labels with bid/ask execution, commissions, pre-expiration exits, and selected-strategy exposure normalization.
  \item Trained one LightGBM regressor using \texttt{strategy} as a categorical feature.
  \item Ran a 2024 test backtest that scores all five strategies, trades at most one strategy per violation, and scales each accepted trade to \$1,000 gross exposure.
\end{itemize}

\section{Engineering Changes Needed For Full Scale}

The first implementation worked on smoke tests but was too slow on the full 2019--2024 universe. The following changes were made to run the full pipeline:

\begin{itemize}
  \item Vectorized TP2/RR2 detection with NumPy instead of nested Python strike-pair loops.
  \item Added parallel detection with \texttt{--n-jobs} and progress logging.
  \item Added the missing per-leg \texttt{spread} feature to detector output.
  \item Replaced row-by-row label calculation with vectorized, chunked label generation.
  \item Made training and backtest load only required columns and valid label rows.
  \item Expanded the JSON backtest summary with exit reasons, daily Sharpe-like ratio, and breakdowns by option type, strategy, and year.
\end{itemize}

\section{Generated Artifacts}

\begin{center}
\begin{tabular}{@{}lrr@{}}
\toprule
Artifact & Rows & Size \\
\midrule
\texttt{data/processed/options\_clean.parquet} & 4,028,173 & 175 MB \\
\texttt{data/features/violations\_features.parquet} & 8,768,014 & 1.4 GB \\
\texttt{data/labels/strategy\_training\_rows.parquet} & 43,840,070 & 6.3 GB \\
\texttt{data/models/lightgbm\_strategy\_return.joblib} & -- & 3.3 MB \\
\texttt{reports/tables/test\_selected\_trades.csv} & 6,132 & 7.1 MB \\
\bottomrule
\end{tabular}
\end{center}

\section{Dataset Expansion Checks}

The first correctness check is that the row counts expand in the expected order: cleaned option quotes, detected violations, then five strategy rows per violation.

\begin{figure}[h]
\centering
\begin{tabular}{@{}lrl@{}}
\toprule
Stage & Rows (millions) & Relative scale \\
\midrule
\barrow{Clean quotes}{4.028}{0.9}{blue}
\barrow{Violations}{8.768}{2.0}{blue}
\barrow{Strategy rows}{43.840}{10.0}{blue}
\bottomrule
\end{tabular}
\caption{Pipeline row-count expansion. The strategy table is exactly five rows per detected violation.}
\end{figure}

\section{Violation And Label Coverage}

\begin{figure}[h]
\centering
\begin{tabular}{@{}lrlrl@{}}
\toprule
Year & Violations (M) & Scale & Valid labels (M) & Scale \\
\midrule
2019 & 1.270 & \textcolor{blue}{\rule{1.6cm}{7pt}} & 3.068 & \textcolor{green!50!black}{\rule{4.2cm}{7pt}}\\
2020 & 1.057 & \textcolor{blue}{\rule{1.3cm}{7pt}} & 2.533 & \textcolor{green!50!black}{\rule{3.5cm}{7pt}}\\
2021 & 1.923 & \textcolor{blue}{\rule{2.4cm}{7pt}} & 4.062 & \textcolor{green!50!black}{\rule{5.5cm}{7pt}}\\
2022 & 0.654 & \textcolor{blue}{\rule{0.8cm}{7pt}} & 1.240 & \textcolor{green!50!black}{\rule{1.7cm}{7pt}}\\
2023 & 1.959 & \textcolor{blue}{\rule{2.4cm}{7pt}} & 3.506 & \textcolor{green!50!black}{\rule{4.7cm}{7pt}}\\
2024 & 1.905 & \textcolor{blue}{\rule{2.4cm}{7pt}} & 3.598 & \textcolor{green!50!black}{\rule{4.9cm}{7pt}}\\
\bottomrule
\end{tabular}
\caption{Detected violations and valid labels by year. Valid labels are strategy-level rows after entry/exit quote checks and pre-expiration handling.}
\end{figure}

\section{Model Training}

The chronological split was preserved:

\begin{center}
\begin{tabular}{@{}ll@{}}
\toprule
Split & Period \\
\midrule
Train & 2019--2022 \\
Validation & 2023 \\
Test & 2024 \\
\bottomrule
\end{tabular}
\end{center}

The full valid training plus validation set was too large for the available memory when fitting the sklearn LightGBM pipeline. To complete a reproducible end-to-end run, the model was trained with explicit caps:

\begin{center}
\begin{tabular}{@{}lr@{}}
\toprule
Metric & Value \\
\midrule
Train rows used & 2,000,000 \\
Validation rows used & 750,000 \\
Validation MAE & 0.14017 \\
Validation RMSE & 0.40161 \\
\bottomrule
\end{tabular}
\end{center}

\section{2024 Test Backtest Results}

\begin{center}
\begin{tabular}{@{}lr@{}}
\toprule
Metric & Value \\
\midrule
Trades & 6,132 \\
Net PnL & \$524,095.82 \\
Average return & 8.55\% \\
Median return & -1.51\% \\
Win rate & 45.58\% \\
Profit factor & 1.6530 \\
Max drawdown & -\$134,245.33 \\
Daily Sharpe-like ratio & 2.31 \\
Average actual holding days & 5.90 \\
\bottomrule
\end{tabular}
\end{center}

The result is positive after scaling each accepted trade to \$1,000 gross exposure. This is still a research backtest, not a validated trading rule; thresholds and portfolio assumptions should be tuned only on validation data.

\begin{figure}[h]
\centering
\begin{tabular}{@{}lrl@{}}
\toprule
Strategy & Trades & Relative scale \\
\midrule
\barrow{\texttt{full4}}{5}{0.1}{blue}
\barrow{\texttt{T1}}{1,690}{7.2}{blue}
\barrow{\texttt{T2}}{247}{1.0}{blue}
\barrow{\texttt{K1}}{2,354}{10.0}{blue}
\barrow{\texttt{K2}}{1,836}{7.8}{blue}
\bottomrule
\end{tabular}
\caption{Accepted 2024 trades by selected strategy. The backtest chooses at most one of the five strategies for each violation.}
\end{figure}

\begin{figure}[h]
\centering
\begin{tabular}{@{}lrl@{}}
\toprule
Month & Fixed-gross PnL & Relative scale \\
\midrule
Jan & -\$31,956 & \textcolor{red}{\rule{1.2cm}{7pt}}\\
Feb & -\$41,695 & \textcolor{red}{\rule{1.6cm}{7pt}}\\
Mar & -\$42,390 & \textcolor{red}{\rule{1.6cm}{7pt}}\\
Apr & \$57,049 & \textcolor{green!50!black}{\rule{2.1cm}{7pt}}\\
May & -\$1,222 & \textcolor{red}{\rule{0.2cm}{7pt}}\\
Jun & -\$32,715 & \textcolor{red}{\rule{1.2cm}{7pt}}\\
Jul & -\$40,662 & \textcolor{red}{\rule{1.5cm}{7pt}}\\
Aug & \$265,971 & \textcolor{green!50!black}{\rule{10.0cm}{7pt}}\\
Sep & \$76,417 & \textcolor{green!50!black}{\rule{2.9cm}{7pt}}\\
Oct & \$33,953 & \textcolor{green!50!black}{\rule{1.3cm}{7pt}}\\
Nov & \$122,264 & \textcolor{green!50!black}{\rule{4.6cm}{7pt}}\\
Dec & \$159,083 & \textcolor{green!50!black}{\rule{6.0cm}{7pt}}\\
\bottomrule
\end{tabular}
\caption{Month-wise fixed-gross PnL for the ML-selected 2024 backtest. Each accepted trade is scaled to \$1,000 gross exposure.}
\end{figure}

\begin{center}
\begin{tabular}{@{}lrrrr@{}}
\toprule
Month & Trades & Monthly PnL & Cumulative PnL & Avg Return \\
\midrule
Jan & 525 & -\$31,956 & -\$31,956 & -6.09\% \\
Feb & 500 & -\$41,695 & -\$73,651 & -8.34\% \\
Mar & 481 & -\$42,390 & -\$116,041 & -8.81\% \\
Apr & 540 & \$57,049 & -\$58,992 & 10.56\% \\
May & 505 & -\$1,222 & -\$60,214 & -0.24\% \\
Jun & 471 & -\$32,715 & -\$92,930 & -6.95\% \\
Jul & 539 & -\$40,662 & -\$133,591 & -7.54\% \\
Aug & 542 & \$265,971 & \$132,379 & 49.07\% \\
Sep & 500 & \$76,417 & \$208,796 & 15.28\% \\
Oct & 575 & \$33,953 & \$242,748 & 5.90\% \\
Nov & 454 & \$122,264 & \$365,012 & 26.93\% \\
Dec & 500 & \$159,083 & \$524,096 & 31.82\% \\
\bottomrule
\end{tabular}
\end{center}

\begin{figure}[h]
\centering
\begin{tabular}{@{}lrl@{}}
\toprule
Date & Cumulative PnL & Loss scale \\
\midrule
Jan 02 & -\$5,125 & \textcolor{red}{\rule{0.4cm}{7pt}}\\
Jan 31 & -\$31,956 & \textcolor{red}{\rule{2.8cm}{7pt}}\\
Feb 29 & -\$73,651 & \textcolor{red}{\rule{6.3cm}{7pt}}\\
Mar 28 & -\$116,041 & \textcolor{red}{\rule{10.0cm}{7pt}}\\
Apr 26 & -\$57,751 & \textcolor{red}{\rule{5.0cm}{7pt}}\\
May 24 & -\$79,452 & \textcolor{red}{\rule{6.8cm}{7pt}}\\
Jun 25 & -\$89,145 & \textcolor{red}{\rule{7.7cm}{7pt}}\\
Jul 24 & -\$3,404 & \textcolor{red}{\rule{0.3cm}{7pt}}\\
Aug 21 & \$38,349 & \textcolor{green!50!black}{\rule{1.1cm}{7pt}}\\
Sep 19 & \$196,587 & \textcolor{green!50!black}{\rule{5.4cm}{7pt}}\\
Oct 17 & \$214,455 & \textcolor{green!50!black}{\rule{5.9cm}{7pt}}\\
Nov 14 & \$367,280 & \textcolor{green!50!black}{\rule{10.0cm}{7pt}}\\
Dec 13 & \$363,317 & \textcolor{green!50!black}{\rule{9.9cm}{7pt}}\\
\bottomrule
\end{tabular}
\caption{Sampled cumulative 2024 PnL path from selected trades after scaling each accepted trade to \$1,000 gross exposure. The final full-year net PnL was \$524,095.82.}
\end{figure}

\section{TP2 Versus RR2 Selection}

In this implementation, TP2 corresponds to call-side violations and RR2 corresponds to put-side violations. The ML selector accepted both types, but they behaved differently after fixed-gross sizing.

\begin{center}
\begin{tabular}{@{}lrrrrrr@{}}
\toprule
Type & Trades & Net PnL & Gross Profit & Gross Loss & Avg Return & Win Rate \\
\midrule
TP2 calls & 733 & \$334,910 & \$371,855 & -\$36,945 & 45.69\% & 59.07\% \\
RR2 puts & 5,399 & \$189,185 & \$954,782 & -\$765,596 & 3.50\% & 43.75\% \\
\bottomrule
\end{tabular}
\end{center}

The TP2 call trades were fewer but cleaner: higher average return, higher win rate, and much smaller gross losses. RR2 put trades contributed positive net PnL too, but with many more trades and much larger offsetting losses. This is useful intuition: the overall 2024 result is not only a high-volume put effect; call-side TP2 trades carried a large share of net profit with fewer selected trades.

\begin{figure}[h]
\centering
\resizebox{\textwidth}{!}{%
\begin{tabular}{@{}lrlrl@{}}
\toprule
Month & TP2 cumulative & Scale & RR2 cumulative & Scale \\
\midrule
Jan & \$24,564 & \textcolor{green!50!black}{\rule{0.7cm}{7pt}} & -\$56,520 & \textcolor{red}{\rule{2.8cm}{7pt}}\\
Feb & \$29,739 & \textcolor{green!50!black}{\rule{0.9cm}{7pt}} & -\$103,391 & \textcolor{red}{\rule{5.1cm}{7pt}}\\
Mar & \$28,791 & \textcolor{green!50!black}{\rule{0.9cm}{7pt}} & -\$144,833 & \textcolor{red}{\rule{7.2cm}{7pt}}\\
Apr & \$32,915 & \textcolor{green!50!black}{\rule{1.0cm}{7pt}} & -\$91,908 & \textcolor{red}{\rule{4.6cm}{7pt}}\\
May & \$58,695 & \textcolor{green!50!black}{\rule{1.8cm}{7pt}} & -\$118,910 & \textcolor{red}{\rule{5.9cm}{7pt}}\\
Jun & \$64,982 & \textcolor{green!50!black}{\rule{1.9cm}{7pt}} & -\$157,911 & \textcolor{red}{\rule{7.8cm}{7pt}}\\
Jul & \$67,867 & \textcolor{green!50!black}{\rule{2.0cm}{7pt}} & -\$201,458 & \textcolor{red}{\rule{10.0cm}{7pt}}\\
Aug & \$173,444 & \textcolor{green!50!black}{\rule{5.2cm}{7pt}} & -\$41,065 & \textcolor{red}{\rule{2.0cm}{7pt}}\\
Sep & \$243,350 & \textcolor{green!50!black}{\rule{7.3cm}{7pt}} & -\$34,554 & \textcolor{red}{\rule{1.7cm}{7pt}}\\
Oct & \$248,272 & \textcolor{green!50!black}{\rule{7.4cm}{7pt}} & -\$5,524 & \textcolor{red}{\rule{0.3cm}{7pt}}\\
Nov & \$332,835 & \textcolor{green!50!black}{\rule{9.9cm}{7pt}} & \$32,178 & \textcolor{green!50!black}{\rule{1.6cm}{7pt}}\\
Dec & \$334,910 & \textcolor{green!50!black}{\rule{10.0cm}{7pt}} & \$189,185 & \textcolor{green!50!black}{\rule{9.4cm}{7pt}}\\
\bottomrule
\end{tabular}
}
\caption{Month-wise cumulative fixed-gross PnL split by TP2 calls and RR2 puts. TP2 calls were positive almost all year; RR2 puts recovered late in the year after a deep mid-year drawdown.}
\end{figure}

\begin{figure}[h]
\centering
\resizebox{\textwidth}{!}{%
\begin{tabular}{@{}lrlrl@{}}
\toprule
Month & TP2 monthly & Scale & RR2 monthly & Scale \\
\midrule
Jan & \$24,564 & \textcolor{green!50!black}{\rule{2.3cm}{7pt}} & -\$56,520 & \textcolor{red}{\rule{3.5cm}{7pt}}\\
Feb & \$5,175 & \textcolor{green!50!black}{\rule{0.5cm}{7pt}} & -\$46,870 & \textcolor{red}{\rule{2.9cm}{7pt}}\\
Mar & -\$948 & \textcolor{red}{\rule{0.2cm}{7pt}} & -\$41,442 & \textcolor{red}{\rule{2.6cm}{7pt}}\\
Apr & \$4,124 & \textcolor{green!50!black}{\rule{0.4cm}{7pt}} & \$52,925 & \textcolor{green!50!black}{\rule{3.3cm}{7pt}}\\
May & \$25,780 & \textcolor{green!50!black}{\rule{2.4cm}{7pt}} & -\$27,002 & \textcolor{red}{\rule{1.7cm}{7pt}}\\
Jun & \$6,286 & \textcolor{green!50!black}{\rule{0.6cm}{7pt}} & -\$39,001 & \textcolor{red}{\rule{2.4cm}{7pt}}\\
Jul & \$2,885 & \textcolor{green!50!black}{\rule{0.3cm}{7pt}} & -\$43,547 & \textcolor{red}{\rule{2.7cm}{7pt}}\\
Aug & \$105,577 & \textcolor{green!50!black}{\rule{10.0cm}{7pt}} & \$160,393 & \textcolor{green!50!black}{\rule{10.0cm}{7pt}}\\
Sep & \$69,906 & \textcolor{green!50!black}{\rule{6.6cm}{7pt}} & \$6,511 & \textcolor{green!50!black}{\rule{0.4cm}{7pt}}\\
Oct & \$4,923 & \textcolor{green!50!black}{\rule{0.5cm}{7pt}} & \$29,030 & \textcolor{green!50!black}{\rule{1.8cm}{7pt}}\\
Nov & \$84,562 & \textcolor{green!50!black}{\rule{8.0cm}{7pt}} & \$37,702 & \textcolor{green!50!black}{\rule{2.4cm}{7pt}}\\
Dec & \$2,076 & \textcolor{green!50!black}{\rule{0.2cm}{7pt}} & \$157,008 & \textcolor{green!50!black}{\rule{9.8cm}{7pt}}\\
\bottomrule
\end{tabular}
}
\caption{Monthly fixed-gross PnL by violation type. RR2 puts drove most of the weak first half and the strong December recovery; TP2 calls were steadier and had smaller loss months.}
\end{figure}

\begin{figure}[h]
\centering
\begin{tabular}{@{}lrl@{}}
\toprule
Exit reason & Trades & Relative scale \\
\midrule
\barrow{5d}{4,209}{10.0}{blue}
\barrow{pre-expiration}{1,923}{4.6}{blue}
\bottomrule
\end{tabular}
\caption{Exit reasons for accepted 2024 trades. Pre-expiration exits confirm the short-DTE handling is active.}
\end{figure}

\section{Interpretation}

The figures support several correctness checks:

\begin{itemize}
  \item The row expansion is mechanically correct: each violation becomes five strategy rows.
  \item Valid labels exist in every train, validation, and test year, so the chronological split is usable.
  \item The backtest traded both calls and puts, but was dominated by puts: 5,399 put trades versus 733 call trades.
  \item The strategy selector used all five strategy labels, though \texttt{K1}, \texttt{K2}, and \texttt{T1} dominated accepted trades.
  \item Pre-expiration exits occurred for 1,923 trades, confirming that labels/backtests do not blindly hold through expiration.
\end{itemize}

The fixed-gross backtest result is positive after costs and portfolio limits, while the unscaled product-replication dollar PnL was not comparable across trades because notionals varied. The implementation is now capable of running the full research loop, but the current thresholds and model are still not a validated trading rule. Next work should focus on validation-only threshold tuning, stronger cost/edge gating, baselines, and memory-aware full-row model training if the uncapped LightGBM fit is required.

\end{document}
